\section{Identification: what one solution determines}\label{f:ch:ident}
\sectionbrief
{$\Omega\subset\R^2$ open and \emph{connected}; $u\in\mathcal D'(\Omega)$; constant coefficients; an operator $P$ with $P(\partial)u=0$ is available (any annihilator, not necessarily the ``true'' one), with irreducible factorisation $P=c\prod_iP_i^{e_i}$.}
{before splitting a solution into pieces one must know its equation. We ask exactly when $u$ alone pins down $P$, and what the intrinsic operator of $u$ is when it does not.}
{first link of the chain in Figure~\ref{fig:pipeline}. Its output is the pair $(G_u,J_u)$, which every later chapter uses, and the test ``$\in\EP$'', which Step~3 of the algorithm implements.}

\subsection{The question in algebraic form}

Suppose $P(\partial)u=0$. Then every multiple $SP$ also kills $u$, so the ideal $(P)$ is contained in $\Ann(u)$:
\[
(P)\subseteq\Ann(u).
\]
We say that $u$ \emph{identifies} $P$ when the two coincide. In that case every equation satisfied by $u$ is a consequence of $P$, and $P$ is determined by $u$ up to a constant factor.

Why can the inclusion be strict? Because $u$ may satisfy extra equations that $P$ knows nothing about. The single Fourier mode of the heat equation is the typical case: it also satisfies $u_{xx}=-\kappa^2u$.

Connectedness of $\Omega$ is necessary. On a disconnected region the two pieces can satisfy different equations, and no criterion on $u$ as a whole can see that.

\subsection{The criterion}

\begin{notebox}{Terminology}
\begin{itemize}
\item $P$ is \emph{reduced} if all $e_i=1$ (no repeated factor).
\item $P$ is \emph{monic in $\tau$} if its leading coefficient as a polynomial in $\tau$ is a nonzero constant. Then $k=\deg_\tau P$ is its order in time, and the \emph{$t$-jet} of $u$ is the vector $(u,\partial_tu,\dots,\partial_t^{k-1}u)$.
\item For each irreducible factor we form the partial quotients
\[
w_i=\Big(\frac{P}{P_i}\Big)(\partial)\,u,\qquad
w_i'=\Big(\frac{P}{P_i^{e_i}}\Big)(\partial)\,u .
\]
In $w_i$ we remove one copy of $P_i$; in $w_i'$ we remove the whole power.
\end{itemize}
\end{notebox}

\begin{theorem}[{Identification criterion \led{C3-T1} \proved}]\label{f:thm:ident}
Let $\Omega$ be connected, $u\in\mathcal D'(\Omega)$ and $P(\partial)u=0$.
\begin{enumerate}[label=(\alph*)]
\item $\Ann(u)=(P)$ if and only if $w_i\notin\EP(\Omega)$ for every irreducible factor $P_i$.
\item $\Var(\Ann u)=\Var(P)$ if and only if $w_i'\notin\EP(\Omega)$ for every $i$.
\item If $P$ is reduced, (a) and (b) coincide.
\item If $P$ is monic in $\tau$, the test can be carried out on the $t$-jet of $u$ (Proposition~\ref{prop:cyclic}).
\end{enumerate}
\dimref{C3-T1}
\end{theorem}

\begin{intuition}
Strip one irreducible factor off $P$ and look at what is left: $w_i$ lies in the kernel of $P_i$.

If $w_i$ is a ``generic'' member of that kernel, not a finite sum of exponentials, then no smaller operator can kill $u$.

If $w_i$ is an exponential-polynomial, then $u$ contains a finite-dimensional piece. That piece is also killed by an operator of lower order, and $u$ cannot tell the large equation from the small one.
\end{intuition}

\begin{figure}[!htb]
\centering
\begin{tikzpicture}
\begin{axis}[width=6.4cm,height=5.2cm,axis lines=middle,xlabel={$\xi$},ylabel={$\tau$},xmin=-2.2,xmax=2.2,ymin=-0.4,ymax=4.6,
 ticks=none,title={\small rich $u$: fills the curve},every axis x label/.style={at={(current axis.right of origin)},anchor=west},
 every axis y label/.style={at={(current axis.above origin)},anchor=south}]
\addplot[cA,line width=3pt,opacity=0.35,domain=-2.1:2.1,samples=60] {x^2};
\addplot[cA,thick,domain=-2.1:2.1,samples=60] {x^2};
\end{axis}
\end{tikzpicture}\hfill
\begin{tikzpicture}
\begin{axis}[width=6.4cm,height=5.2cm,axis lines=middle,xlabel={$\xi$},ylabel={$\tau$},xmin=-2.2,xmax=2.2,ymin=-0.4,ymax=4.6,
 ticks=none,title={\small poor $u$: three points},every axis x label/.style={at={(current axis.right of origin)},anchor=west},
 every axis y label/.style={at={(current axis.above origin)},anchor=south}]
\addplot[cG,thick,domain=-2.1:2.1,samples=60] {x^2};
\addplot[only marks,mark=*,cB,mark size=2.8pt] coordinates {(1,1) (-1,1) (2,4)};
\end{axis}
\end{tikzpicture}
\caption{The heat curve $\tau=\xi^2$. Left: a solution that is not an exponential-polynomial (the heat kernel) uses the whole curve, and $\Ann(u)=(P)$. Right: a finite sum of exponentials lives on finitely many points. Then $\Ann(u)$ contains polynomials that vanish at those points but are not multiples of $P$.}\label{fig:rich-poor}
\end{figure}

\begin{proposition}[{Uniqueness of the identified operator}]\label{f:prop:unique}
If $\Ann(u)=(P)$, then every operator $Q$ with $Q(\partial)u=0$ is a multiple $Q=SP$ with $S\in\ring$. In particular $P$ is, up to a nonzero constant, the only annihilator of minimal degree.

\dimref{C3-T1}
\end{proposition}

The proof is one line: $Q\in(P)$ means $Q=SP$, and degrees add, so $\deg Q\ge\deg P$ with equality only when $S$ is a constant.

\begin{example}[{The heat equation, three cases}]\label{ex:heat3}
Take $P=\tau-D\xi^2$ with $D>0$. It is irreducible, so $P/P_1=1$ and the criterion reads
\[
\Ann(u)=(P)\iff u\notin\EP .
\]

\textbf{One Fourier mode.} Let $u=e^{-D\kappa^2t+i\kappa x}$. Then $u_t=-D\kappa^2u$ and $u_{xx}=-\kappa^2u$. So $u$ is also killed by $\partial_x-i\kappa$ and by $\partial_t+D\kappa^2$:
\[
\Ann(u)=(\tau+D\kappa^2,\ \xi-i\kappa).
\]
This is the ideal of one point of $\Var(P)$, strictly larger than $(P)$. Indeed $P$ belongs to it:
\[
P=(\tau+D\kappa^2)-D\,(\xi-i\kappa)(\xi+i\kappa).
\]

\textbf{A finite superposition.} For $u=\sum_{j=1}^rc_je^{-D\kappa_j^2t+i\kappa_jx}$ with distinct $\kappa_j$, $u\in\EP$. Its annihilator is the ideal of the $r$ points $(-D\kappa_j^2,\,i\kappa_j)$. It contains many operators that are not multiples of $P$.

\textbf{The heat kernel.} Let $u=t^{-1/2}e^{-x^2/(4Dt)}$ for $t>0$. Its $x$-derivatives are polynomials of growing degree times the same Gaussian. They span an infinite-dimensional space, so $u\notin\EP$ (Lemma~\ref{f:lem:ep}). Hence $\Ann(u)=(P)$: every constant-coefficient equation satisfied by the heat kernel is a multiple of $\tau-D\xi^2$.

\emph{Message.} A ``poor'' solution remembers only finitely many exponentials and cannot tell the heat equation from many others. A ``rich'' solution remembers the whole equation.
\end{example}

\begin{example}[{The wave equation and d'Alembert \led{C3-T5} \proved}]\label{ex:wave-ident}
Let $P=(\tau-c\xi)(\tau+c\xi)$ with real $c\ne0$, and $u=f(x-ct)+g(x+ct)$. The two partial quotients are
\begin{align*}
w_1&=(\partial_t+c\partial_x)\,u=2c\,g'(x+ct),\\
w_2&=(\partial_t-c\partial_x)\,u=-2c\,f'(x-ct).
\end{align*}
So the criterion says:
\[
\Ann(u)=(P)\iff f\notin\EP\ \text{ and }\ g\notin\EP .
\]
A single travelling wave ($g=0$) never identifies the wave equation: it also satisfies the first-order transport equation. The same theorem covers Klein--Gordon $\tau^2-\xi^2+\mu^2$ and the telegraph operator $\tau^2+\alpha\tau-c^2\xi^2$: both are irreducible conics, so they are identified iff $u\notin\EP$. For the beam $\tau^2+\xi^4=(\tau-i\xi^2)(\tau+i\xi^2)$ both Schr\"odinger parts must be outside $\EP$.

\dimref{C3-T5}
\end{example}

\begin{takeaway}
Theorem~\ref{f:thm:ident} turns ``does $u$ identify $P$?'' into finitely many membership tests ``$\in\EP$''. Each test is a finite-dimensional computation (Lemma~\ref{f:lem:ep}). Whether the tests can be carried out from the data one actually has is the subject of Section~\ref{f:ch:data}.
\end{takeaway}

\subsection{The intrinsic operator and the spectral form}\label{sec:spectral}

Theorem~\ref{f:thm:ident} says when $(P)$ is the whole annihilator. When it is not, what does $u$ determine? The answer uses the two pieces of $\Ann(u)$ given by Lemma~\ref{f:lem:gcd}.

\begin{notebox}{Point spectrum}
Recall $\Ann(u)=G_u\,J_u$ with $G_u=\Gcd\Ann(u)$. The \emph{point spectrum} of $u$ is the finite set
\[
\sigma_p(u)=\Var(J_u)\subset\C^2 .
\]
Its points are the exponents $(\tau_0,\xi_0)$ of the exponentials $e^{\tau_0t+\xi_0x}$ that are ``present'' in $u$. The word spectrum is meant in the algebraic sense, not as the spectrum of an operator on a Hilbert space. We write $G_u\sim P$ when the two agree up to a nonzero constant.
\end{notebox}

\begin{theorem}[{Spectral characterisation \led{P4-T2} \proved}]\label{thm:spec}
Let $P$ be an annihilator of $u$. Then
\[
\Ann(u)=(P)\iff \sigma_p(u)=\emptyset\ \text{ and }\ G_u\sim P .
\]
Moreover the exponential-polynomial part of the derivative space is the finite quotient
\[
\ring u\cap\EP\;\cong\;\ring/J_u .
\]
\dimref{P4-T2}
\end{theorem}

\begin{example}[{Reading the spectrum}]
\textbf{One exponential.} For $u=e^{\tau_0t+\xi_0x}$: $\Ann(u)=(\tau-\tau_0,\ \xi-\xi_0)$, so $G_u=1$ and $\sigma_p(u)=\{(\tau_0,\xi_0)\}$. A single exponential has no infinite part and a one-point spectrum.

\textbf{Two exponentials.} For $u=e^{\tau_1t+\xi_1x}+e^{\tau_2t+\xi_2x}$ with distinct exponents: $G_u=1$ and $\sigma_p(u)$ has two points.

\textbf{A rich solution.} If $\Ann(u)=(P)$ then $\sigma_p(u)=\emptyset$ and $G_u\sim P$.
\end{example}

The operator $G_u$ is the correct ``operator intrinsic to $u$''. It depends only on $u$, while $P$ depends on which annihilator was chosen. The next theorem computes it from any annihilator.

\begin{theorem}[{Intrinsic operator \led{C3-T6} \proved}]\label{thm:intrinsic}
Let $P=c\prod_iP_i^{e_i}$ be any annihilator of $u$. For each $i$ let
\[
k_i=\max\Big\{k\le e_i:\ \Big(\frac{P}{P_i^{k}}\Big)(\partial)u\in\EP\Big\}
\]
(with $k_i=0$ if no such $k\ge1$ exists), and $f_i=e_i-k_i$. Then
\[
G_u\sim\prod_iP_i^{f_i},
\]
and $\Ann(u)$ is principal if and only if $G_u(\partial)u=0$. Computing $G_u$ takes at most $\sum_ie_i$ tests ``$\in\EP$''.

\dimref{C3-T6}
\end{theorem}

\begin{example}[{A repeated factor}]
Let $G$ be the heat kernel ($D=1$), $P=(\tau-\xi^2)^3$ and $u=t\,G$. Then:
\begin{align*}
\Big(\frac{P}{P_1}\Big)(\partial)u&=(\partial_t-\partial_x^2)^2(tG)=0\in\EP,\\
\Big(\frac{P}{P_1^2}\Big)(\partial)u&=(\partial_t-\partial_x^2)(tG)=G\notin\EP .
\end{align*}
So $k_1=1$, $f_1=2$, and
\[
G_u=(\tau-\xi^2)^2,\qquad \Ann(u)=(G_u).
\]
The function $tG$ identifies the \emph{square} of the heat operator, not its cube.
\end{example}

\begin{caution}[title={Why $G_u$ and not the variety}]
A tempting stronger statement replaces $G_u$ by the variety $\Var(\Ann u)$. It is correct only for reduced $P$, because zero sets do not see multiplicities.

\emph{Example.} Let $P=\tau^2$ and $u=e^{x^2}$. Then $u_{tt}=0$. No nonzero polynomial in $\partial_x$ alone kills $u$, so
\[
\Ann(u)=(\tau).
\]
Here $\sigma_p(u)=\emptyset$, but $G_u=\tau\not\sim P$. So $\Ann(u)\ne(P)$, while $\Var(\Ann u)=\{\tau=0\}=\Var(P)$. The variety alone would wrongly say that $u$ identifies $\tau^2$.
\end{caution}

\subsection{Refinements}

\subsubsection{Evolutions: identification inside a finite class}

For an evolution $u_t=a(\partial_x)u$ the symbol $P=\tau-a(\xi)$ has degree one in $\tau$, so it is irreducible. The criterion says: $\Ann(u)=(P)$ iff $u\notin\EP$. When $u$ \emph{is} an exponential-polynomial, we can still ask whether $a$ is recovered among polynomials of degree at most $n$.

\begin{theorem}[{Evolutions \led{C3-T4} \proved}]\label{thm:evol}
Let $P=\tau-a(\xi)$ and $P(\partial)u=0$ on a connected $\Omega$.
\begin{enumerate}[label=(\alph*)]
\item $\Ann(u)=(P)$ if and only if $u\notin\EP(\Omega)$.
\item If $u=\sum_kq_k\,e^{\tau_kt+\xi_kx}$ with distinct exponents and nonzero $q_k$, put $\delta_k=\deg_xq_k$ and
\[
D=\sum_k(\delta_k+1).
\]
If $\deg a\le n$, then among all $b$ of degree at most $n$, the evolutions $\tau-b$ that kill $u$ are exactly
\[
b\in a+\prod_k(\xi-\xi_k)^{\delta_k+1}\cdot\C[\xi]_{\le n-D}.
\]
So $a$ is identified in degree $\le n$ if and only if $D\ge n+1$.
\end{enumerate}
\dimref{C3-T4}
\end{theorem}

This is Hermite interpolation: a polynomial of degree $n$ is fixed by $n+1$ conditions, and each mode $q_k\,e^{\tau_kt+\xi_kx}$ supplies $\delta_k+1$ of them.

\begin{example}[{Advection--diffusion from two modes}]
Take $u_t=-cu_x+Du_{xx}$, so $a(\xi)=-c\xi+D\xi^2$ has degree $n=2$. Two simple modes give $D=2<3$: the pair $(c,D)$ is not identified, because one free parameter remains. Three distinct modes give $D=3$: both coefficients are fixed.
\end{example}

\subsubsection{A naive transfer that fails}

On a periodic domain an earlier result said: if the Fourier expansion of $u$ has infinitely many frequencies on every factor curve, then $\Ann(u)=(P)$ \led{R4-T5}. On a \emph{bounded} domain this transfer is false, because the expansion is not unique.

\begin{falsebox}{The spectrum of an expansion is not intrinsic \led{C3-C1} \falsest}\label{cex:C3C1}
Let $P=\tau+\xi$ and $\Omega=(0,1)^2$. On $(-\pi,\pi)$ the Fourier series
\[
e^{s}=\sum_{k\in\Z}c_k\,e^{iks},\qquad c_k=\frac{\sinh(\pi)(-1)^k}{\pi(1-ik)}\neq0,
\]
converges in $L^2$. So $u=e^{x-t}=\sum_kc_k\,e^{ik(x-t)}$ has infinitely many exponents on the line $\tau+\xi=0$. Yet
\[
\Ann(u)=(\tau+1,\ \xi-1)\supsetneq(\tau+\xi).
\]
\dimref{C3-C1}
\end{falsebox}

The criterion must therefore be stated on $u$ itself, as in Theorem~\ref{f:thm:ident}, not on an expansion of $u$.

\subsubsection{Symmetries and local data}

\begin{proposition}[{Translation symmetry \led{C3-P3} \pw}]\label{prop:sym}
Suppose $u$ is invariant under translations along a direction, $\ell(\partial)u=0$ with $\ell=\alpha\tau+\beta\xi\ne0$, and $P(\partial)u=0$ with $P$ nonconstant. If $P\not\sim\ell$, then $\Ann(u)\ne(P)$. If moreover $\ell\nmid P$, then $u\in\EP$.

\dimref{C3-P3}
\end{proposition}

So a symmetry with constant characteristic destroys identification, and it is a special case of the criterion. Symmetries with variable coefficients are invisible in $\C[\partial]$: the heat kernel is scale invariant, $(2t\partial_t+x\partial_x+1)G=0$, yet $\Ann(G)=(\tau-\xi^2)$. Identifiability depends on the ring of coefficients.

\begin{proposition}[{Local data \led{C3-P5} \pw}]\label{prop:local}
Identification can be lost on a smaller region. Let $P=\tau+\xi$, $F(s)=e^s+\varphi(s)$ with $\varphi$ a smooth bump supported in $(2,3)$, and $u=F(x-t)$. On $(0,1)\times(-5,5)$, $u\notin\EP$ and $\Ann(u)=(\tau+\xi)$. On the smaller square $(0,1)^2$, $u=e^{x-t}$ and the annihilator is the ideal of a point.

For elliptic $P$, every solution is real-analytic, and then $\Ann(u|_{\Omega'})=\Ann(u)$ for every connected $\Omega'\subset\Omega$.

\dimref{C3-P5}
\end{proposition}

\begin{keypoint}
A solution identifies its equation exactly when, after removing any one irreducible factor, what is left is not a finite sum of exponentials. In general, $u$ identifies only its intrinsic operator $G_u$. The difference between $P$ and $G_u$ is a finite set of exponentials, the point spectrum $\sigma_p(u)$.
\end{keypoint}

